% ═══════════════════════════════════════════════════════════════════════════
% REFERÊNCIAS
% ═══════════════════════════════════════════════════════════════════════════
\begin{thebibliography}{99}

\bibitem{flavell1979}
FLAVELL, J. H. Metacognition and cognitive monitoring. \textit{American Psychologist}, v. 34, n. 10, p. 906--911, 1979. \doi{10.1037/0003-066X.34.10.906}.

\bibitem{cox2005}
COX, M. T. Metacognition in computation. \textit{Artificial Intelligence}, v. 169, n. 2, p. 104--141, 2005. \doi{10.1016/j.artint.2005.10.009}.

\bibitem{minsky2006}
MINSKY, M. \textit{The Emotion Machine}. New York: Simon \& Schuster, 2006. ISBN: 978-0-7432-7663-8.

\bibitem{baars1988}
BAARS, B. J. \textit{A Cognitive Theory of Consciousness}. Cambridge: Cambridge University Press, 1988. \doi{10.1017/CBO9780511551326}.

\bibitem{baars1997}
BAARS, B. J. \textit{In the Theater of Consciousness}. Oxford: Oxford University Press, 1997. \doi{10.1093/acprof:oso/9780195102659.001.1}.

\bibitem{graziano2013}
GRAZIANO, M. S. A. \textit{Consciousness and the Social Brain}. Oxford: Oxford University Press, 2013. ISBN: 978-0-19-992864-4.

\bibitem{graziano2019}
GRAZIANO, M. S. A. \textit{Rethinking Consciousness}. New York: W. W. Norton, 2019. ISBN: 978-0-393-65261-1.

\bibitem{atkinson1968}
ATKINSON, R. C.; SHIFFRIN, R. M. Human memory. In: SPENCE, K. W. (Ed.). \textit{The Psychology of Learning and Motivation}. Academic Press, 1968. v. 2, p. 89--195. \doi{10.1016/S0079-7421(08)60422-3}.

\bibitem{ostrom1990}
OSTROM, E. \textit{Governing the Commons}. Cambridge: Cambridge University Press, 1990. \doi{10.1017/CBO9780511807763}.

\bibitem{ostrom2010}
OSTROM, E. Beyond markets and states. \textit{American Economic Review}, v. 100, n. 3, p. 641--672, 2010. \doi{10.1257/aer.100.3.641}.

\bibitem{ostrom1999}
OSTROM, E. et al. Revisiting the commons. \textit{Science}, v. 284, n. 5412, p. 278--282, 1999. \doi{10.1126/science.284.5412.278}.

\bibitem{smith1973}
SMITH, J. M.; PRICE, G. R. The logic of animal conflict. \textit{Nature}, v. 246, p. 15--18, 1973. \doi{10.1038/246015a0}.

\bibitem{axelrod1984}
AXELROD, R. \textit{The Evolution of Cooperation}. New York: Basic Books, 1984. ISBN: 978-0-465-02121-5.

\bibitem{beck2003}
BECK, K. \textit{Test-Driven Development: By Example}. Boston: Addison-Wesley, 2003. ISBN: 978-0-321-14653-3.

\bibitem{gamma1994}
GAMMA, E. et al. \textit{Design Patterns}. Reading: Addison-Wesley, 1994. ISBN: 978-0-201-63361-0.

\bibitem{polya1945}
PÓLYA, G. \textit{How to Solve It}. Princeton: Princeton University Press, 1945. \doi{10.2307/j.ctt7swnh}.

\bibitem{lakatos1976}
LAKATOS, I. \textit{Proofs and Refutations}. Cambridge: Cambridge University Press, 1976. \doi{10.1017/CBO9781139171472}.

\bibitem{granger1969}
GRANGER, C. W. J. Investigating causal relations. \textit{Econometrica}, v. 37, n. 3, p. 424--438, 1969. \doi{10.2307/1912791}.

\bibitem{miller1956}
MILLER, G. A. The magical number seven. \textit{Psychological Review}, v. 63, n. 2, p. 81--97, 1956. \doi{10.1037/h0043158}.

\bibitem{hoare1969}
HOARE, C. A. R. An axiomatic basis for computer programming. \textit{Communications of the ACM}, v. 12, n. 10, p. 576--580, 1969. \doi{10.1145/363235.363259}.

\bibitem{dijkstra1976}
DIJKSTRA, E. W. \textit{A Discipline of Programming}. Prentice-Hall, 1976. ISBN: 978-0-13-215871-8.

\bibitem{gentzen1935}
GENTZEN, G. Untersuchungen über das logische Schließen. \textit{Mathematische Zeitschrift}, v. 39, p. 176--210, 1935. \doi{10.1007/BF01201353}.

\bibitem{kahneman2011}
KAHNEMAN, D. \textit{Thinking, Fast and Slow}. New York: Farrar, Straus and Giroux, 2011. ISBN: 978-0-374-27563-1.

\bibitem{russell2010}
RUSSELL, S.; NORVIG, P. \textit{Artificial Intelligence: A Modern Approach}. 3. ed. Prentice Hall, 2010. ISBN: 978-0-13-604259-4.

\bibitem{clarke1999}
CLARKE, E. M.; GRUMBERG, O.; PELED, D. A. \textit{Model Checking}. MIT Press, 1999. ISBN: 978-0-262-03270-4.

\bibitem{knuth1968}
KNUTH, D. E. \textit{The Art of Computer Programming, Vol. 1}. Addison-Wesley, 1968. ISBN: 978-0-201-03801-9.

\bibitem{peffers2007}
PEFFERS, K. et al. A design science research methodology. \textit{Journal of Management Information Systems}, v. 24, n. 3, p. 45--77, 2007. \doi{10.2753/MIS0742-1222240302}.

\bibitem{hevner2004}
HEVNER, A. R. et al. Design science in information systems research. \textit{MIS Quarterly}, v. 28, n. 1, p. 75--105, 2004. \doi{10.2307/25148625}.

\bibitem{janzen2005}
JANZEN, D.; SAIEDIAN, H. Test-driven development: Concepts, taxonomy, and future direction. \textit{Computer}, v. 38, n. 9, p. 43--50, 2005. \doi{10.1109/MC.2005.314}.

\bibitem{schmidhuber2007}
SCHMIDHUBER, J. Gödel machines: Fully self-referential optimal general problem solvers. In: GOERTZEL, B.; PENNACHIN, C. (Eds.). \textit{Artificial General Intelligence}. Springer, 2007. p. 199--226.

\bibitem{wooldridge2009}
WOOLDRIDGE, M. \textit{An Introduction to MultiAgent Systems}. 2. ed. Wiley, 2009. ISBN: 978-0-470-51946-2.

\bibitem{bonabeau1999}
BONABEAU, E.; DORIGO, M.; THERAULAZ, G. \textit{Swarm Intelligence}. Oxford University Press, 1999. ISBN: 978-0-19-513159-8.

\bibitem{nexo2026}
WAZIONAPPS. NEXO Brain. \textit{GitHub}, 2026. \url{https://github.com/wazionapps/nexo}. Acesso: 10 jun. 2026.

\bibitem{autoresearch2026}
ALVINREAL. Awesome Autoresearch. \textit{GitHub}, 2026. \url{https://github.com/alvinreal/awesome-autoresearch}. Acesso: 10 jun. 2026.

\bibitem{ruflo2026}
RUVNET. Ruflo. \textit{GitHub}, 2026. \url{https://github.com/ruvnet/ruflo}. Acesso: 10 jun. 2026.

\bibitem{opencode2026}
LARANJEIRA, M. C. OpenCode Ecosystem v5.4.0 R23. \textit{GitHub}, 2026. \url{https://github.com/MarceloClaro/OpenCode_Ecosystem}. Acesso: 10 jun. 2026.

\bibitem{planning2026}
ADIO, O. Planning-with-files. \textit{GitHub}, 2026. \url{https://github.com/OthmanAdi/planning-with-files}. Acesso: 10 jun. 2026.

\end{thebibliography}
